% Figure: three buoyancy regimes side by side: sinking (denser than fluid),
% rising (less dense, fully submerged), and floating (weight balances
% buoyancy with partial submersion). No em-dashes.
\begin{figure}[htbp]
\centering
\begin{tikzpicture}[scale=1.0]
% three panels share a waterline
\foreach \x in {0,4,8}{
  \fill[engblue!12] (\x,0) rectangle (\x+3,2.4);
  \draw[engblue, thick] (\x,2.4) -- (\x+3,2.4);
}
% Panel 1: sinking object near bottom (denser)
\fill[engblack!70] (1.2,0.2) rectangle (1.8,0.8);
\draw[engred, ->] (1.5,0.8) -- (1.5,1.5); % buoyancy up
\node[engred, font=\scriptsize, anchor=west] at (1.5,1.2) {$F_B$};
\draw[engblack, ->] (1.5,0.8) -- (1.5,0.1); % weight down (longer)
\node[engblack, font=\scriptsize, anchor=west] at (1.5,0.35) {$W$};
\node[font=\small, anchor=south] at (1.5,2.4) {sinks};
\node[font=\scriptsize, anchor=north] at (1.5,0.0) {$\rho_b>\rho_f$};
% Panel 2: rising object (less dense, fully submerged)
\fill[englime!70] (5.2,1.0) rectangle (5.8,1.6);
\draw[engred, ->] (5.5,1.6) -- (5.5,2.3); % buoyancy up (longer)
\node[engred, font=\scriptsize, anchor=west] at (5.5,2.0) {$F_B$};
\draw[engblack, ->] (5.5,1.6) -- (5.5,1.1); % weight down (shorter)
\node[engblack, font=\scriptsize, anchor=west] at (5.5,1.25) {$W$};
\node[font=\small, anchor=south] at (5.5,2.4) {rises};
\node[font=\scriptsize, anchor=north] at (5.5,0.0) {$\rho_b<\rho_f$};
% Panel 3: floating object at the surface (partial submersion)
\fill[engamber!70] (9.2,2.0) rectangle (9.8,2.9);
\draw[engblue, thick] (8,2.4) -- (11,2.4);
\draw[engred, ->] (9.5,2.0) -- (9.5,1.3); % buoyancy (drawn down for room) equal
\node[engred, font=\scriptsize, anchor=east] at (9.5,1.65) {$F_B$};
\draw[engblack, ->] (9.5,2.9) -- (9.5,3.5);
\node[engblack, font=\scriptsize, anchor=west] at (9.5,3.2) {$W=F_B$};
\node[font=\small, anchor=south] at (9.5,2.95) {floats};
\node[font=\scriptsize, anchor=north] at (9.5,0.0) {$V_{\mathrm{disp}}=m/\rho_f$};
\end{tikzpicture}
\caption[The three buoyancy regimes]{Archimedes' principle in three
cases. Left: a body denser than the fluid has weight exceeding the
maximum buoyancy and sinks. Centre: a fully submerged body less dense
than the fluid has buoyancy exceeding its weight and rises to the
surface. Right: a floating body submerges just enough that the displaced
volume's weight equals its own, $V_{\mathrm{disp}}=m/\rho_f$. In every
case the buoyancy force is the weight of the displaced fluid,
$F_B=\rho_fgV_{\mathrm{disp}}$, acting upward through the centroid of the
displaced volume.}
\label{fig:vol03ch10:buoyancy-regimes}
\end{figure}
